\section{相关工作}\label{sec:related}

\paragraph{多相机 BEV 感知}
鸟瞰（bird's-eye-view, BEV）表征将多相机图像中的视觉证据映射到机器人或车辆坐标系的俯视平面，使不同视角的观测能够在统一空间网格中进行融合、推理和规划。
以 BEV 为中间表征的视觉感知方法通过预定义空间查询把多相机图像特征对齐到该统一坐标系。
显式提升范式将图像特征沿深度分布投影至 BEV，代表性工作包括 Lift--Splat--Shoot、BEVDet、
BEVDepth 和 BEVStereo\cite{philion2020lift,huang2021bevdet,li2022bevdepth,li2022bevstereo}；
MatrixVT 则进一步讨论了高效的多相机到 BEV 变换\cite{zhou2022matrixvt}。
查询式方法则利用三维参考点、位置嵌入或跨视图变换直接读取相机特征\cite{wang2022detr3d,li2022petr,liu2023petrv2,li2023fbbev}，
并可在共享 BEV 中融合多模态信息\cite{liu2023bevfusion}。
BEVFormer 进一步以空间交叉注意力和时序自注意力建立 BEV 查询之间的几何关联\cite{li2022bevformer}。
这些方法为统一坐标系下的视觉感知奠定基础。

\paragraph{三维占用预测与高效表征}
视觉三维占用将自由空间和语义实体编码为稠密或压缩的体素场。MonoScene、TPVFormer 和 VoxFormer 分别从单目特征提升、
三平面表征和可见稀疏体素出发完成场景补全\cite{cao2022monoscene,huang2023tpvformer,li2023voxformer}；
SurroundOcc 与 OpenOccupancy 推动了多相机稠密占用预测和环视基准的发展\cite{wei2023surroundocc,wang2023openoccupancy}。
由于三维体素的计算代价随空间分辨率快速增长，后续工作从局部--全局结构、
紧凑或由粗到细的查询、稀疏交互，以及高效解码和渲染监督等方向降低开销\cite{zhang2023occformer,ma2024cotr,wang2024panoocc,liu2024sparseocc,tang2024sparseocc,yu2023flashocc,pan2024renderocc}。
占用表征也进一步扩展为时序世界模型：BEVDet4D、OccWorld、IR-WM 与 ForecastOcc 分别
利用跨帧 BEV、离散场景 token 或历史图像观测预测动态环境和未来状态\cite{huang2022bevdet4d,zheng2024occworld,mei2025irwm,mohan2026forecastocc}。
这些研究验证了占用作为统一空间接口的价值，但其普遍面向全向交通环境，
依赖更广的相机覆盖或历史运动信息。

\paragraph{农业与越野三维感知数据}
城市三维感知的常用数据集包括 KITTI\cite{geiger2012kitti}、SemanticKITTI\cite{behley2019semantickitti} 和 nuScenes\cite{caesar2020nuscenes}，
但它们的传感器布置、类别构成和道路结构与农田不同。
农业场景中的公开数据多面向作物行检测、二维语义分割、实例分割或特定作业目标，难以同时提供多相机标定和稠密三维语义体素；
越野方向的 RELLIS-3D\cite{jiang2020rellis} 也主要服务于多模态分割与建图。
虚拟驾驶平台 CARLA\cite{dosovitskiy2017carla} 及领域随机化\cite{tobin2017domainrandomization,peng2018sim2real}、
语义一致域适应\cite{hoffman2018cycada} 说明了仿真环境可以极大的缓和真实世界数据收集成本高、
耗时久的痛点，尤其是在农业场景。
在真实越野方向，ORAD-3D 提供涵盖农田、林地、草地和道路等场景的三维占用
基准\cite{chen2025orad3d},为本文提供了真实场景下的验证途径。
